\documentclass[12pt]{article}
\usepackage[utf8]{inputenc}
\usepackage[T1]{fontenc}
\usepackage{hyperref}
\title{Sediq Language Resources for K-12 Curriculum Development}
\author{Compiled by Claude Sonnet 5 (claude-sonnet-5) for Andr\'as Kornai}
\date{}
\begin{document}
\maketitle
\section{Introduction}
Sediq (Seediq; ISO 639-3 trv) is an Austronesian language of Taiwan with three main dialects: Truku, Toda (Tuuda) and Tgdaya (Tkdaya, Paran). Searches turned up little material suited to a K-12 curriculum. This dossier lists only what the searches returned.
\section{Primary Educational Resources}
No textbook series or Ministry of Education curriculum documents were found in the search results. Omniglot reports that a Latin-alphabet orthography exists and that Seediq has been taught in some schools since 2001; this is secondary information. \url{https://www.omniglot.com/writing/seediq.htm}
\section{Linguistic and Reference Materials}
\begin{itemize}
\item Overview of the language: \url{https://en.wikipedia.org/wiki/Seediq_language}
\item A Reference Grammar of Seediq by Henry Yung-li Chang (2000, Formosan Languages Series No.~6, Taipei: Yuanliou). It is cited in a National Tsing Hua University course syllabus: \url{https://khub.nthu.edu.tw/document/courses/09820GPTS530100.pdf}
\item Blog overview of Seediq, which notes that language certification tests are divided by dialect: \url{https://thelanguagecloset.com/2020/12/12/languages-of-taiwan-seediq-kari-sediq-kari-seediq-kari-seejiq/}
\end{itemize}
\section{Religious and Community Resources}
None were identified in the search results.
\section{Digital and Audio Resources}
No dictionaries, corpora or audio were found in the search results. Searches in Chinese (for example the terms for Seediq dictionary and teaching materials) and a check of Taiwan's Ministry of Education indigenous-language resources would be the next step.
\section{Curriculum Development Context}
Materials appear to be organised by dialect (Truku, Toda, Tgdaya), so a curriculum will need dialect-specific materials. The sources found are thin and mostly secondary.
\end{document}